\documentclass[11pt,a4paper]{article}

\usepackage[utf8]{inputenc}
\usepackage[T1]{fontenc}
\usepackage[french]{babel}
\usepackage{amsmath,amssymb}
\usepackage{geometry}
\usepackage{enumitem}
\usepackage{multicol}
\usepackage{tabularx}
\usepackage{booktabs}
\usepackage{xcolor}
\usepackage{mdframed}
\usepackage{lmodern}
\usepackage{microtype}
\usepackage{fancyhdr}
\usepackage{lastpage}
\usepackage{parskip}

\geometry{top=2cm, bottom=2.5cm, left=2.2cm, right=2.2cm}

\pagestyle{fancy}
\fancyhf{}
\fancyhead[L]{\small\textit{Mathématiques - Corrigé pédagogique}}
\fancyhead[R]{\small\textit{Fonctions de plusieurs variables \& Opérateurs vectoriels}}
\fancyfoot[C]{\small Page \thepage\ / \pageref{LastPage}}
\renewcommand{\headrulewidth}{0.3pt}
\renewcommand{\footrulewidth}{0pt}

% Couleurs
\definecolor{grisbox}{gray}{0.94}
\definecolor{bleutexte}{RGB}{30,60,120}
\definecolor{vertok}{RGB}{0,120,60}
\definecolor{rougeerr}{RGB}{180,30,30}
\definecolor{orangeinfo}{RGB}{200,100,0}
\definecolor{bleuclair}{RGB}{220,230,245}
\definecolor{vertclair}{RGB}{220,245,230}
\definecolor{orangeclair}{RGB}{255,240,220}

% Environnements colorés
\newmdenv[backgroundcolor=vertclair, linecolor=vertok, linewidth=1pt,
          innerleftmargin=10pt, innerrightmargin=10pt,
          innertopmargin=6pt, innerbottommargin=6pt,
          skipabove=4pt, skipbelow=6pt]{reponsebox}

\newmdenv[backgroundcolor=bleuclair, linecolor=bleutexte, linewidth=1pt,
          innerleftmargin=10pt, innerrightmargin=10pt,
          innertopmargin=6pt, innerbottommargin=6pt,
          skipabove=4pt, skipbelow=4pt]{methodebox}

\newmdenv[backgroundcolor=orangeclair, linecolor=orangeinfo, linewidth=1pt,
          innerleftmargin=10pt, innerrightmargin=10pt,
          innertopmargin=6pt, innerbottommargin=6pt,
          skipabove=4pt, skipbelow=6pt]{piegebox}

\newmdenv[backgroundcolor=grisbox, linewidth=0pt,
          innerleftmargin=10pt, innerrightmargin=10pt,
          innertopmargin=6pt, innerbottommargin=6pt,
          skipabove=6pt, skipbelow=6pt]{notebox}

% Style des sections
\usepackage{titlesec}
\titleformat{\section}[block]{\large\bfseries\color{bleutexte}}{}{0pt}{}[\vspace{-4pt}\rule{\linewidth}{0.4pt}]
\titlespacing{\section}{0pt}{16pt}{8pt}

% En-tête de question
\newcounter{qnum}
\newcommand{\qtitre}[2]{%
  \stepcounter{qnum}%
  \medskip\noindent
  \textbf{\color{bleutexte}Question \theqnum}\enspace
  \textit{(#1)}\enspace
  \textcolor{vertok}{\textbf{$\Rightarrow$ Réponse : #2}}
  \par\smallskip
}

% Énoncé de la question
\newcommand{\qenonce}[1]{%
  \begin{mdframed}[backgroundcolor=grisbox, linewidth=0pt,
    innerleftmargin=8pt, innerrightmargin=8pt,
    innertopmargin=5pt, innerbottommargin=5pt,
    skipabove=2pt, skipbelow=4pt]
  \small\textit{#1}
  \end{mdframed}
}

\begin{document}

% ─── TITRE ───────────────────────────────────────────────────────────────────
\begin{center}
  {\Large\bfseries\color{bleutexte} Corrigé TOMIC Blanc}\\[4pt]
  {\large Contrôle - Fonctions de plusieurs variables \& Opérateurs vectoriels}\\[6pt]
  {\small\itshape
    Pour chaque question : la réponse correcte, la méthode complète et les pièges à éviter.}
\end{center}

\vspace{2pt}
\noindent\rule{\linewidth}{0.6pt}
\vspace{4pt}

% ── Grille de correction ──────────────────────────────────────────────────────
\begin{notebox}
\textbf{Grille de correction rapide}\\[4pt]
\begin{tabular}{@{}lll@{\quad}lll@{\quad}lll@{}}
Q1  & QCM    & \textbf{B}                          &
Q11 & QCM    & \textbf{B}                          &
Q21 & QCM    & \textbf{B}                          \\
Q2  & QCM    & \textbf{C}                          &
Q12 & QCM    & \textbf{C}                          &
Q22 & QCM    & \textbf{B}                          \\
Q3  & Calcul & limite inexistante                  &
Q13 & QCM    & \textbf{C}                          &
Q23 & QCM    & \textbf{B}                          \\
Q4  & QCM    & \textbf{C}                          &
Q14 & QCM    & \textbf{B}                          &
Q24 & QCM    & \textbf{C}                          \\
Q5  & QCM    & \textbf{C}                          &
Q15 & Calcul & $(1,2)$                             &
Q25 & Calcul & $3x+z\;;\;(-y,-z,0)$               \\
Q6  & QCM    & \textbf{B}                          &
Q16 & QCM    & \textbf{B}                          &
Q26 & QCM    & \textbf{C}                          \\
Q7  & Calcul & $z=2x+y-1$                         &
Q17 & Calcul & min en $(1,2)$                      &
Q27 & QCM    & \textbf{B}                          \\
Q8  & QCM    & \textbf{A}                          &
Q18 & QCM    & \textbf{C}                          &
Q28 & Calcul & $x+2y+2z=9$                        \\
Q9  & QCM    & \textbf{B}                          &
Q19 & QCM    & \textbf{B}                          \\
Q10 & Calcul & $3x^2\cos(y)$                       &
Q20 & QCM    & \textbf{B}                          \\
\end{tabular}
\end{notebox}

% ═══════════════════════════════════════════════════════════════════════════════
\section{Partie 1 - Limites et continuité}
% ═══════════════════════════════════════════════════════════════════════════════

%────── Q1 ────────────────────────────────────────────────────────────────────
\qtitre{QCM}{B - $\lim_{(x,y)\to a}f(x,y) = f(a)$}
\qenonce{Soit $f:\mathbb{R}^2\to\mathbb{R}$. Elle est continue en $a=(a_1,a_2)$ si et seulement si :
\quad A)~Elle admet des dérivées partielles en $a$
\quad B)~$\lim_{(x,y)\to a}f(x,y) = f(a)$
\quad C)~Elle est bornée au voisinage de $a$
\quad D)~$f(a) = 0$}

\begin{methodebox}
\textbf{Rappel.} La continuité en $a$ est définie par l'égalité entre la limite et la valeur :
$$f \text{ continue en } a \;\Longleftrightarrow\; \lim_{(x,y)\to a} f(x,y) = f(a)$$
\end{methodebox}

\begin{itemize}[leftmargin=1.5em, itemsep=2pt]
  \item \textcolor{rougeerr}{\textbf{A)}} L'existence des dérivées partielles ne suffit \emph{pas} à garantir la continuité : contre-exemple classique avec $f(x,y)=\frac{xy}{x^2+y^2}$.
  \item \textcolor{vertok}{\textbf{B)}} Définition exacte de la continuité. \textbf{Correct.}
  \item \textcolor{rougeerr}{\textbf{C)}} Être borné n'implique pas la continuité.
  \item \textcolor{rougeerr}{\textbf{D)}} Annuler $f$ en $a$ n'a rien à voir avec la continuité.
\end{itemize}

%────── Q2 ────────────────────────────────────────────────────────────────────
\qtitre{QCM}{C - $F(x,y,z) = 0$}
\qenonce{Une surface de $\mathbb{R}^3$ donnée sous forme implicite s'écrit :
\quad A)~$z = f(x,y)$
\quad B)~$\mathbf{r}(t) = (x(t),y(t),z(t))$
\quad C)~$F(x,y,z) = 0$
\quad D)~$(u,v)\mapsto(r_1,r_2,r_3)$}

\begin{methodebox}
\textbf{Trois représentations d'une surface :}
\begin{itemize}[itemsep=1pt, topsep=2pt]
  \item \textbf{Explicite :} $z = f(x,y)$ - la surface est le graphe de $f$.
  \item \textbf{Implicite :} $F(x,y,z) = 0$ - ensemble de niveau d'une fonction de 3 variables.
  \item \textbf{Paramétrique :} $(u,v)\mapsto\mathbf{r}(u,v)$ - deux paramètres décrivent la surface.
\end{itemize}
La forme B, $\mathbf{r}(t)$, est un paramétrage de \emph{courbe} (une variable).
\end{methodebox}

%────── Q3 (Calcul) ───────────────────────────────────────────────────────────
\qtitre{Calcul}{La limite n'existe pas}
\qenonce{Soit $f(x,y) = \dfrac{xy}{x^2+y^2}$ pour $(x,y)\neq(0,0)$. Montrer que la limite en $(0,0)$ n'existe pas en testant les chemins $y=0$ et $y=x$.}

\begin{reponsebox}
\textbf{Principe.} Si deux chemins vers $(0,0)$ donnent des limites différentes, la limite globale n'existe pas.

\textbf{Chemin 1 - axe $Ox$ ($y = 0$) :}
$$f(x, 0) = \frac{x \cdot 0}{x^2 + 0} = 0 \xrightarrow[x\to 0]{} \mathbf{0}$$

\textbf{Chemin 2 - diagonale ($y = x$) :}
$$f(x, x) = \frac{x \cdot x}{x^2 + x^2} = \frac{x^2}{2x^2} = \frac{1}{2} \xrightarrow[x\to 0]{} \mathbf{\dfrac{1}{2}}$$

Les deux limites $0 \neq \tfrac{1}{2}$ sont différentes, donc :
$$\boxed{\lim_{(x,y)\to(0,0)}\frac{xy}{x^2+y^2} \text{ n'existe pas.}}$$
\end{reponsebox}

\begin{piegebox}
\textbf{Piège.} Tester uniquement $y=0$ donne 0, et on pourrait conclure que la limite est 0. Il faut toujours tester plusieurs chemins - la non-existence se prouve par deux chemins donnant des valeurs distinctes.
\end{piegebox}

% ═══════════════════════════════════════════════════════════════════════════════
\section{Partie 2 - Différentielle et plan tangent}
% ═══════════════════════════════════════════════════════════════════════════════

%────── Q4 ────────────────────────────────────────────────────────────────────
\qtitre{QCM}{C - $df = 2xy\,dx + (x^2+\cos y)\,dy$}
\qenonce{La différentielle de $f(x,y) = x^2 y + \sin(y)$ est :
\quad A)~$2y\,dx + x^2\,dy$
\quad B)~$2xy\,dx + x^2\,dy$
\quad C)~$2xy\,dx + (x^2+\cos y)\,dy$
\quad D)~$(2xy+x^2)\,d(x+y)$}

\begin{methodebox}
\textbf{Différentielle.} $df = \dfrac{\partial f}{\partial x}dx + \dfrac{\partial f}{\partial y}dy$
\end{methodebox}

Calcul des dérivées partielles :
$$\frac{\partial f}{\partial x} = 2xy \qquad (\text{on dérive } x^2y \text{ par rapport à } x \text{ ; }\sin y \text{ est une constante})$$
$$\frac{\partial f}{\partial y} = x^2 + \cos y \qquad (\text{on dérive } x^2y \text{ et } \sin y \text{ par rapport à } y)$$

\textbf{Pourquoi B est faux :} oubli du terme $\cos y$ dans $\dfrac{\partial(\sin y)}{\partial y}$.

%────── Q5 ────────────────────────────────────────────────────────────────────
\qtitre{QCM}{C - Continue en $a$}
\qenonce{Si $f$ est différentiable en $a$, alors $f$ est nécessairement :
\quad A)~De classe $\mathcal{C}^1$ au voisinage de $a$
\quad B)~Deux fois dérivable en $a$
\quad C)~Continue en $a$
\quad D)~Analytique en $a$}

\begin{methodebox}
\textbf{Hiérarchie des régularités (du plus fort au plus faible) :}
$$\mathcal{C}^\infty \;\Rightarrow\; \mathcal{C}^2 \;\Rightarrow\; \mathcal{C}^1 \;\Rightarrow\; \text{différentiable} \;\Rightarrow\; \text{continue} \;\Rightarrow\; \text{dérivées partielles}$$
La différentiabilité implique la continuité, mais pas l'inverse.
\end{methodebox}

\begin{itemize}[leftmargin=1.5em, itemsep=2pt]
  \item \textcolor{rougeerr}{\textbf{A)}} $\mathcal{C}^1$ implique différentiabilité, mais pas l'inverse.
  \item \textcolor{rougeerr}{\textbf{B)}} La différentiabilité ne garantit pas la dérivabilité au 2\textsuperscript{e} ordre.
  \item \textcolor{vertok}{\textbf{C)}} Différentiable $\Rightarrow$ continue. \textbf{Correct.}
  \item \textcolor{rougeerr}{\textbf{D)}} L'analyticité est bien plus forte (développement en série entière).
\end{itemize}

%────── Q6 ────────────────────────────────────────────────────────────────────
\qtitre{QCM}{B - $z = f(a,b) + \frac{\partial f}{\partial x}(a,b)(x-a) + \frac{\partial f}{\partial y}(a,b)(y-b)$}
\qenonce{L'équation du plan tangent à $z=f(x,y)$ au point $(a,b,f(a,b))$ est :
\quad A)~$z = f(a,b)\cdot(x+y)$
\quad B)~$z = f(a,b) + \frac{\partial f}{\partial x}(a,b)(x-a) + \frac{\partial f}{\partial y}(a,b)(y-b)$
\quad C)~$z = \frac{\partial f}{\partial x}(a,b)\,x + \frac{\partial f}{\partial y}(a,b)\,y$
\quad D)~$z = f'(a)(x-a)$}

\begin{methodebox}
\textbf{Plan tangent = développement limité d'ordre 1.} Le plan tangent en $(a,b)$ est l'approximation affine :
$$z = f(a,b) + \frac{\partial f}{\partial x}(a,b)(x-a) + \frac{\partial f}{\partial y}(a,b)(y-b)$$
C'est aussi l'équation du plan d'équation $\nabla F(a,b,f(a,b))\cdot(M-M_0)=0$ avec $F(x,y,z) = z - f(x,y)$.
\end{methodebox}

La proposition C oublie le terme $f(a,b)$ (n'est valable qu'au voisinage de l'origine si $f(0,0)=0$). La proposition D est une formule en une seule variable.

%────── Q7 (Calcul) ───────────────────────────────────────────────────────────
\qtitre{Calcul}{$df = 2xe^y\,dx + x^2e^y\,dy$ \quad et \quad plan tangent : $z = 2x+y-1$}
\qenonce{Soit $f(x,y) = x^2 e^y$. Calculer la différentielle $df$, puis l'équation du plan tangent en $(1,0)$.}

\begin{reponsebox}
\textbf{Calcul des dérivées partielles :}
$$\frac{\partial f}{\partial x} = 2xe^y \qquad\qquad \frac{\partial f}{\partial y} = x^2e^y$$

\textbf{Différentielle :}
$$\boxed{df = 2xe^y\,dx + x^2e^y\,dy}$$

\textbf{Valeurs au point $(1,0)$ :}
$$f(1,0) = 1^2\cdot e^0 = 1,\qquad \frac{\partial f}{\partial x}(1,0) = 2\cdot 1\cdot e^0 = 2,\qquad \frac{\partial f}{\partial y}(1,0) = 1^2\cdot e^0 = 1$$

\textbf{Équation du plan tangent :}
$$z = 1 + 2(x-1) + 1(y-0) = 1 + 2x - 2 + y$$
$$\boxed{z = 2x + y - 1}$$

\textbf{Vérification :} En $(1,0)$ : $2(1)+0-1 = 1 = f(1,0)$ \checkmark
\end{reponsebox}

% ═══════════════════════════════════════════════════════════════════════════════
\section{Partie 3 - Dérivées partielles et théorème de Schwartz}
% ═══════════════════════════════════════════════════════════════════════════════

%────── Q8 ────────────────────────────────────────────────────────────────────
\qtitre{QCM}{A - $\dfrac{\partial f}{\partial x} = g'(h)\cdot\dfrac{\partial h}{\partial x}$}
\qenonce{Si $f(x,y)=g(h(x,y))$, la règle de la chaîne donne :
\quad A)~$\frac{\partial f}{\partial x} = g'(h)\cdot\frac{\partial h}{\partial x}$
\quad B)~$\frac{\partial f}{\partial x} = g'(h) + \frac{\partial h}{\partial x}$
\quad C)~$\frac{\partial f}{\partial x} = g(h)\cdot h(x,y)$
\quad D)~$\frac{\partial f}{\partial x} = g''(h)$}

\begin{methodebox}
\textbf{Règle de la chaîne (composition).} Si $f = g \circ h$ avec $h : \mathbb{R}^2 \to \mathbb{R}$ et $g : \mathbb{R} \to \mathbb{R}$ :
$$\frac{\partial f}{\partial x} = g'(h(x,y))\cdot\frac{\partial h}{\partial x}, \qquad \frac{\partial f}{\partial y} = g'(h(x,y))\cdot\frac{\partial h}{\partial y}$$
\end{methodebox}

\textbf{Exemple :} $f(x,y) = \sin(x^2 y)$, $h = x^2 y$, $g = \sin$.
$$\frac{\partial f}{\partial x} = \cos(x^2y)\cdot 2xy \qquad \frac{\partial f}{\partial y} = \cos(x^2y)\cdot x^2$$

%────── Q9 ────────────────────────────────────────────────────────────────────
\qtitre{QCM}{B - $\dfrac{\partial^2 f}{\partial x\partial y} = \dfrac{\partial^2 f}{\partial y\partial x}$}
\qenonce{Le théorème de Schwartz affirme que, pour $f\in\mathcal{C}^2$ :
\quad A)~$\frac{\partial^2 f}{\partial x^2} = \frac{\partial^2 f}{\partial y^2}$
\quad B)~$\frac{\partial^2 f}{\partial x\partial y} = \frac{\partial^2 f}{\partial y\partial x}$
\quad C)~$\frac{\partial f}{\partial x} = \frac{\partial f}{\partial y}$
\quad D)~$\Delta f = 0$}

\begin{methodebox}
\textbf{Théorème de Clairaut-Schwartz.} Pour $f\in\mathcal{C}^2$ sur un ouvert, l'ordre de dérivation n'a pas d'importance :
$$\frac{\partial^2 f}{\partial x\partial y} = \frac{\partial^2 f}{\partial y\partial x}$$
\textbf{Conséquence pratique :} pour calculer $\dfrac{\partial^2 f}{\partial x\partial y}$, on peut dériver dans l'ordre qui est le plus simple.
\end{methodebox}

\begin{itemize}[leftmargin=1.5em, itemsep=2pt]
  \item \textcolor{rougeerr}{\textbf{A)}} Faux en général : $\frac{\partial^2 f}{\partial x^2} = \frac{\partial^2 f}{\partial y^2}$ n'est pas une propriété générale.
  \item \textcolor{vertok}{\textbf{B)}} Exactement l'énoncé du théorème. \textbf{Correct.}
  \item \textcolor{rougeerr}{\textbf{C, D)}} Sans rapport avec Schwartz.
\end{itemize}

%────── Q10 (Calcul) ──────────────────────────────────────────────────────────
\qtitre{Calcul}{$\dfrac{\partial^2 f}{\partial x\partial y} = \dfrac{\partial^2 f}{\partial y\partial x} = 3x^2\cos(y)$}
\qenonce{Soit $f(x,y)=x^3\sin(y)$. Calculer $\frac{\partial f}{\partial x}$, $\frac{\partial f}{\partial y}$, $\frac{\partial^2 f}{\partial x^2}$ et $\frac{\partial^2 f}{\partial x\partial y}$. Vérifier Schwartz.}

\begin{reponsebox}
\textbf{Dérivées partielles d'ordre 1 :}
$$\frac{\partial f}{\partial x} = 3x^2\sin(y), \qquad \frac{\partial f}{\partial y} = x^3\cos(y)$$

\textbf{Dérivées partielles d'ordre 2 :}
$$\frac{\partial^2 f}{\partial x^2} = \frac{\partial}{\partial x}\bigl(3x^2\sin y\bigr) = 6x\sin(y)$$
$$\frac{\partial^2 f}{\partial x\partial y} = \frac{\partial}{\partial y}\bigl(3x^2\sin y\bigr) = \boxed{3x^2\cos(y)}$$

\textbf{Vérification du théorème de Schwartz :}
$$\frac{\partial^2 f}{\partial y\partial x} = \frac{\partial}{\partial x}\bigl(x^3\cos y\bigr) = 3x^2\cos(y) = \frac{\partial^2 f}{\partial x\partial y} \quad\checkmark$$
\end{reponsebox}

% ═══════════════════════════════════════════════════════════════════════════════
\section{Partie 4 - Équations aux dérivées partielles}
% ═══════════════════════════════════════════════════════════════════════════════

%────── Q11 ────────────────────────────────────────────────────────────────────
\qtitre{QCM}{B - Du premier ordre, linéaire homogène}
\qenonce{L'équation $\dfrac{\partial u}{\partial x} + 2\dfrac{\partial u}{\partial y} = 0$ est :
\quad A)~Du second ordre, linéaire
\quad B)~Du premier ordre, linéaire homogène
\quad C)~Du premier ordre, non linéaire
\quad D)~Une équation des ondes}

\begin{methodebox}
\textbf{Classification des EDP.}
\begin{itemize}[itemsep=1pt, topsep=2pt]
  \item \textbf{Ordre :} ordre de la dérivée partielle la plus élevée.
  \item \textbf{Linéaire :} $u$ et toutes ses dérivées apparaissent à la puissance 1, sans produits entre elles.
  \item \textbf{Homogène :} second membre nul.
\end{itemize}
\end{methodebox}

Ici : seules des dérivées d'ordre 1 ($\frac{\partial u}{\partial x}$, $\frac{\partial u}{\partial y}$) $\Rightarrow$ \textbf{premier ordre}. Dépendance linéaire en $u_x$ et $u_y$ $\Rightarrow$ \textbf{linéaire}. Second membre $= 0$ $\Rightarrow$ \textbf{homogène}.

%────── Q12 ────────────────────────────────────────────────────────────────────
\qtitre{QCM}{C - $\dfrac{\partial^2 u}{\partial t^2} = c^2\dfrac{\partial^2 u}{\partial x^2}$}
\qenonce{L'équation de propagation des ondes en une dimension s'écrit :
\quad A)~$\frac{\partial u}{\partial t} = k\frac{\partial^2 u}{\partial x^2}$
\quad B)~$\frac{\partial^2 u}{\partial x^2}+\frac{\partial^2 u}{\partial y^2}=0$
\quad C)~$\frac{\partial^2 u}{\partial t^2}=c^2\frac{\partial^2 u}{\partial x^2}$
\quad D)~$\frac{\partial u}{\partial t}+u\frac{\partial u}{\partial x}=0$}

\begin{methodebox}
\textbf{Les trois grandes EDP classiques :}
\begin{center}
\begin{tabular}{@{}lll@{}}
\toprule
Équation & Forme & Type \\
\midrule
Ondes (propagation) & $u_{tt} = c^2 u_{xx}$ & Hyperbolique \\
Chaleur (diffusion) & $u_t = k\,u_{xx}$ & Parabolique \\
Laplace (stationnaire) & $u_{xx}+u_{yy}=0$ & Elliptique \\
\bottomrule
\end{tabular}
\end{center}
\end{methodebox}

La proposition D est l'équation de Burgers (non linéaire). La proposition B est l'équation de Laplace.

%────── Q13 ────────────────────────────────────────────────────────────────────
\qtitre{QCM}{C - $\xi = x+ct,\;\eta = x-ct$}
\qenonce{Pour résoudre l'équation des ondes par la méthode de d'Alembert, le changement de variables est :
\quad A)~$\xi=x^2$, $\eta=t^2$
\quad B)~$\xi=x/c$, $\eta=t/c$
\quad C)~$\xi=x+ct$, $\eta=x-ct$
\quad D)~$\xi=\cos(x)$, $\eta=\sin(ct)$}

\begin{methodebox}
\textbf{Méthode de d'Alembert.} Le changement $\xi = x+ct$, $\eta = x-ct$ transforme :
$$\frac{\partial^2 u}{\partial t^2} = c^2\frac{\partial^2 u}{\partial x^2} \quad\longrightarrow\quad \frac{\partial^2 u}{\partial\xi\,\partial\eta} = 0$$
La solution générale est $u(x,t) = f(x+ct) + g(x-ct)$ : superposition d'une onde progressive et d'une onde régressive.
\end{methodebox}

\begin{piegebox}
\textbf{Interprétation physique.} $f(x-ct)$ se déplace vers la droite à la vitesse $c$, et $g(x+ct)$ vers la gauche. Tout signal se propage sans déformation (contrairement à l'équation de la chaleur).
\end{piegebox}

% ═══════════════════════════════════════════════════════════════════════════════
\section{Partie 5 - Extrema d'une fonction de deux variables}
% ═══════════════════════════════════════════════════════════════════════════════

%────── Q14 ────────────────────────────────────────────────────────────────────
\qtitre{QCM}{B - $\nabla f(a,b) = \mathbf{0}$}
\qenonce{Un point $(a,b)$ est un point critique de $f$ si et seulement si :
\quad A)~$f(a,b) = 0$
\quad B)~$\nabla f(a,b) = \mathbf{0}$
\quad C)~La hessienne de $f$ en $(a,b)$ est nulle
\quad D)~$f$ n'est pas définie en $(a,b)$}

\begin{methodebox}
\textbf{Point critique (ou point stationnaire).} $(a,b)$ est un point critique si :
$$\nabla f(a,b) = \mathbf{0} \;\Longleftrightarrow\; \frac{\partial f}{\partial x}(a,b) = 0 \;\text{ et }\; \frac{\partial f}{\partial y}(a,b) = 0$$
\textbf{Condition nécessaire} pour un extremum local (si $f$ est différentiable). La hessienne permet ensuite de \emph{classifier} le point critique.
\end{methodebox}

%────── Q15 (Calcul) ──────────────────────────────────────────────────────────
\qtitre{Calcul}{Point critique : $(1, 2)$}
\qenonce{Trouver les points critiques de $f(x,y) = x^2 + y^2 - 2x - 4y + 1$.}

\begin{reponsebox}
\textbf{Calcul du gradient et résolution de $\nabla f = \mathbf{0}$ :}
$$\frac{\partial f}{\partial x} = 2x - 2 = 0 \quad\Rightarrow\quad x = 1$$
$$\frac{\partial f}{\partial y} = 2y - 4 = 0 \quad\Rightarrow\quad y = 2$$

$$\boxed{\text{Point critique unique : }(1, 2)}$$

\textbf{Valeur en ce point :} $f(1,2) = 1 + 4 - 2 - 8 + 1 = -4$.
\end{reponsebox}

%────── Q16 ────────────────────────────────────────────────────────────────────
\qtitre{QCM}{B - Un minimum local}
\qenonce{Si $H = \frac{\partial^2 f}{\partial x^2}\frac{\partial^2 f}{\partial y^2} - \bigl(\frac{\partial^2 f}{\partial x\partial y}\bigr)^2 > 0$ et $\frac{\partial^2 f}{\partial x^2}(a,b) > 0$, alors $(a,b)$ est :
\quad A)~Un maximum local
\quad B)~Un minimum local
\quad C)~Un point selle
\quad D)~On ne peut pas conclure}

\begin{methodebox}
\textbf{Critère du discriminant hessien $H$.} En un point critique $(a,b)$ de $f\in\mathcal{C}^2$ :
\begin{center}
\begin{tabular}{@{}lll@{}}
\toprule
Condition & Nature & \\
\midrule
$H > 0$ et $f_{xx} > 0$ & Minimum local & paraboloïde vers le haut \\
$H > 0$ et $f_{xx} < 0$ & Maximum local & paraboloïde vers le bas \\
$H < 0$                 & Point selle   & \\
$H = 0$                 & Cas douteux   & test insuffisant \\
\bottomrule
\end{tabular}
\end{center}
\end{methodebox}

%────── Q17 (Calcul) ──────────────────────────────────────────────────────────
\qtitre{Calcul}{$(1,2)$ est un minimum local, $f(1,2) = -4$}
\qenonce{Pour la fonction de la question précédente, calculer $H$ au point critique $(1,2)$ et conclure.}

\begin{reponsebox}
\textbf{Dérivées partielles secondes de $f(x,y) = x^2+y^2-2x-4y+1$ :}
$$\frac{\partial^2 f}{\partial x^2} = 2, \qquad \frac{\partial^2 f}{\partial y^2} = 2, \qquad \frac{\partial^2 f}{\partial x\partial y} = 0$$

\textbf{Discriminant hessien :}
$$H = \frac{\partial^2 f}{\partial x^2}\cdot\frac{\partial^2 f}{\partial y^2} - \left(\frac{\partial^2 f}{\partial x\partial y}\right)^2 = 2\times 2 - 0^2 = 4 > 0$$

\textbf{Conclusion :} $H = 4 > 0$ et $\dfrac{\partial^2 f}{\partial x^2} = 2 > 0$, donc :
$$\boxed{(1,2) \text{ est un \textbf{minimum local}, avec } f(1,2) = -4}$$

\textbf{Interprétation :} $f(x,y) = (x-1)^2 + (y-2)^2 - 4$ est un paraboloïde de révolution de sommet $(1,2,-4)$.
\end{reponsebox}

%────── Q18 ────────────────────────────────────────────────────────────────────
\qtitre{QCM}{C - $H < 0$}
\qenonce{Un point col (point selle) est un point critique $(a,b)$ tel que :
\quad A)~$H>0$ et $f_{xx}>0$
\quad B)~$H>0$ et $f_{xx}<0$
\quad C)~$H < 0$
\quad D)~$H = 0$}

\begin{methodebox}
\textbf{Point selle.} Si $H < 0$, la fonction \emph{croît} dans certaines directions et \emph{décroît} dans d'autres depuis $(a,b)$. Ce n'est ni un maximum ni un minimum : c'est un col, comme la surface d'une selle de cheval.
\end{methodebox}

\textbf{Exemple :} $f(x,y) = x^2 - y^2$ en $(0,0)$ : $\nabla f(0,0) = \mathbf{0}$, $H = -4 < 0$ $\Rightarrow$ point selle.

% ═══════════════════════════════════════════════════════════════════════════════
\section{Partie 6 - Opérateurs vectoriels}
% ═══════════════════════════════════════════════════════════════════════════════

%────── Q19 ────────────────────────────────────────────────────────────────────
\qtitre{QCM}{B - $\nabla f = \bigl(\frac{\partial f}{\partial x},\,\frac{\partial f}{\partial y},\,\frac{\partial f}{\partial z}\bigr)$}
\qenonce{Le gradient d'un champ scalaire $f(x,y,z)$ est :
\quad A)~$\frac{\partial f}{\partial x}+\frac{\partial f}{\partial y}+\frac{\partial f}{\partial z}$
\quad B)~$\bigl(\frac{\partial f}{\partial x},\,\frac{\partial f}{\partial y},\,\frac{\partial f}{\partial z}\bigr)$
\quad C)~$\frac{\partial^2 f}{\partial x^2}+\frac{\partial^2 f}{\partial y^2}+\frac{\partial^2 f}{\partial z^2}$
\quad D)~$f\cdot\bigl(\frac{\partial}{\partial x}+\frac{\partial}{\partial y}+\frac{\partial}{\partial z}\bigr)$}

\begin{methodebox}
\textbf{Gradient.} $\nabla f = \mathrm{grad}\,f = \left(\dfrac{\partial f}{\partial x},\;\dfrac{\partial f}{\partial y},\;\dfrac{\partial f}{\partial z}\right)$ est un \textbf{vecteur}.

La proposition A est un scalaire (somme, pas vecteur). La proposition C est le laplacien $\Delta f$.
\end{methodebox}

%────── Q20 ────────────────────────────────────────────────────────────────────
\qtitre{QCM}{B - Perpendiculaire aux surfaces de niveau, dans le sens croissant}
\qenonce{En tout point, le vecteur gradient $\nabla f$ est orienté :
\quad A)~Parallèlement aux surfaces de niveau de $f$
\quad B)~Perpendiculairement aux surfaces de niveau, dans le sens des valeurs croissantes
\quad C)~Dans la direction de l'axe $Oz$
\quad D)~Vers le minimum de $f$}

\begin{methodebox}
\textbf{Interprétation géométrique du gradient.}
\begin{itemize}[itemsep=1pt, topsep=2pt]
  \item $\nabla f$ est \textbf{perpendiculaire} aux surfaces de niveau $\{f = c\}$.
  \item $\nabla f$ pointe dans la \textbf{direction de plus forte croissance} de $f$.
  \item $\|\nabla f\|$ mesure le taux de variation maximal.
  \item En particulier : $\nabla f = \mathbf{0}$ aux points critiques (extrema, selles).
\end{itemize}
\end{methodebox}

%────── Q21 ────────────────────────────────────────────────────────────────────
\qtitre{QCM}{B - $\text{rot}\,\mathbf{F} = \bigl(\frac{\partial R}{\partial y}-\frac{\partial Q}{\partial z},\;\frac{\partial P}{\partial z}-\frac{\partial R}{\partial x},\;\frac{\partial Q}{\partial x}-\frac{\partial P}{\partial y}\bigr)$}
\qenonce{Le rotationnel de $\mathbf{F}=(P,Q,R)$ est :
\quad A)~$\frac{\partial P}{\partial x}+\frac{\partial Q}{\partial y}+\frac{\partial R}{\partial z}$
\quad B)~$\bigl(\frac{\partial R}{\partial y}-\frac{\partial Q}{\partial z},\;\frac{\partial P}{\partial z}-\frac{\partial R}{\partial x},\;\frac{\partial Q}{\partial x}-\frac{\partial P}{\partial y}\bigr)$
\quad C)~$\bigl(\frac{\partial P}{\partial x},\;\frac{\partial Q}{\partial y},\;\frac{\partial R}{\partial z}\bigr)$
\quad D)~$\nabla\cdot\mathbf{F}$}

\begin{methodebox}
\textbf{Rotationnel.} $\text{rot}\,\mathbf{F} = \nabla\times\mathbf{F}$ est un \textbf{vecteur} :
$$\text{rot}\,\mathbf{F} = \begin{vmatrix}\mathbf{i}&\mathbf{j}&\mathbf{k}\\\partial_x&\partial_y&\partial_z\\P&Q&R\end{vmatrix}
= \left(\frac{\partial R}{\partial y}-\frac{\partial Q}{\partial z},\;\frac{\partial P}{\partial z}-\frac{\partial R}{\partial x},\;\frac{\partial Q}{\partial x}-\frac{\partial P}{\partial y}\right)$$
La proposition A est la \textbf{divergence} (scalaire). La proposition D $\nabla\cdot\mathbf{F}$ est aussi la divergence.
\end{methodebox}

%────── Q22 ────────────────────────────────────────────────────────────────────
\qtitre{QCM}{B - $\text{rot}\,\mathbf{F} = \mathbf{0}$}
\qenonce{Si $\mathbf{F}=\nabla f$ est un champ gradient, alors :
\quad A)~$\text{div}\,\mathbf{F}=0$
\quad B)~$\text{rot}\,\mathbf{F}=\mathbf{0}$
\quad C)~$\Delta f = 0$
\quad D)~$\mathbf{F}=\mathbf{0}$}

\begin{methodebox}
\textbf{Propriété fondamentale :} $\text{rot}(\nabla f) = \mathbf{0}$ pour tout $f\in\mathcal{C}^2$.

\textbf{Preuve (composante $z$) :}
$$\frac{\partial}{\partial x}\frac{\partial f}{\partial y} - \frac{\partial}{\partial y}\frac{\partial f}{\partial x} = \frac{\partial^2 f}{\partial x\partial y} - \frac{\partial^2 f}{\partial y\partial x} = 0 \quad \text{(Schwartz)}$$

Un champ dont le rotationnel est nul est dit \textbf{irrotationnel} ou \textbf{conservatif}. Tout champ gradient est conservatif.
\end{methodebox}

\begin{piegebox}
\textbf{Réciproque.} Sur un domaine simplement connexe, un champ irrotationnel ($\text{rot}\,\mathbf{F}=\mathbf{0}$) est un gradient. Sur un domaine non simplement connexe, ce n'est plus vrai en général.
\end{piegebox}

%────── Q23 ────────────────────────────────────────────────────────────────────
\qtitre{QCM}{B - $\text{div}\,\mathbf{F} = \dfrac{\partial P}{\partial x}+\dfrac{\partial Q}{\partial y}+\dfrac{\partial R}{\partial z}$}
\qenonce{La divergence de $\mathbf{F}=(P,Q,R)$ est :
\quad A)~$\bigl(\frac{\partial Q}{\partial z}-\frac{\partial R}{\partial y},\;\ldots\bigr)$
\quad B)~$\frac{\partial P}{\partial x}+\frac{\partial Q}{\partial y}+\frac{\partial R}{\partial z}$
\quad C)~$\sqrt{P^2+Q^2+R^2}$
\quad D)~$\frac{\partial P}{\partial y}+\frac{\partial Q}{\partial z}+\frac{\partial R}{\partial x}$}

\begin{methodebox}
\textbf{Divergence.} $\text{div}\,\mathbf{F} = \nabla\cdot\mathbf{F} = \dfrac{\partial P}{\partial x}+\dfrac{\partial Q}{\partial y}+\dfrac{\partial R}{\partial z}$ est un \textbf{scalaire}.

\textbf{Interprétation physique :} mesure le taux d'expansion local du champ (flux sortant par unité de volume). Un champ à divergence nulle est dit \textbf{solénoïdal} (incompressible en mécanique des fluides).
\end{methodebox}

La proposition A est le rotationnel. La proposition C est la norme de $\mathbf{F}$.

%────── Q24 ────────────────────────────────────────────────────────────────────
\qtitre{QCM}{C - $\Delta f = \dfrac{\partial^2 f}{\partial x^2}+\dfrac{\partial^2 f}{\partial y^2}+\dfrac{\partial^2 f}{\partial z^2}$}
\qenonce{L'opérateur laplacien d'un champ scalaire $f$ est :
\quad A)~$\bigl(\frac{\partial f}{\partial x},\frac{\partial f}{\partial y},\frac{\partial f}{\partial z}\bigr)$
\quad B)~$\frac{\partial f}{\partial x}\cdot\frac{\partial f}{\partial y}\cdot\frac{\partial f}{\partial z}$
\quad C)~$\frac{\partial^2 f}{\partial x^2}+\frac{\partial^2 f}{\partial y^2}+\frac{\partial^2 f}{\partial z^2}$
\quad D)~$\text{rot}(\nabla f)$}

\begin{methodebox}
\textbf{Laplacien.} $\Delta f = \text{div}(\nabla f) = \dfrac{\partial^2 f}{\partial x^2}+\dfrac{\partial^2 f}{\partial y^2}+\dfrac{\partial^2 f}{\partial z^2}$ est un \textbf{scalaire}.

La proposition A est le gradient (vecteur). La proposition D est nulle par Schwartz ($\text{rot}(\nabla f) = \mathbf{0}$).
\end{methodebox}

%────── Q25 (Calcul) ──────────────────────────────────────────────────────────
\qtitre{Calcul}{$\text{div}\,\mathbf{F} = 3x+z$ \quad et \quad $\text{rot}\,\mathbf{F} = (-y,\,-z,\,0)$}
\qenonce{Soit $\mathbf{F}(x,y,z) = (x^2,\;yz,\;xz)$. Calculer $\text{div}\,\mathbf{F}$ et $\text{rot}\,\mathbf{F}$.}

\begin{reponsebox}
\textbf{Divergence} ($P = x^2$, $Q = yz$, $R = xz$) :
$$\text{div}\,\mathbf{F} = \frac{\partial(x^2)}{\partial x}+\frac{\partial(yz)}{\partial y}+\frac{\partial(xz)}{\partial z} = 2x + z + x = \boxed{3x+z}$$

\textbf{Rotationnel :}
$$\text{rot}\,\mathbf{F} = \left(\frac{\partial(xz)}{\partial y}-\frac{\partial(yz)}{\partial z},\;\frac{\partial(x^2)}{\partial z}-\frac{\partial(xz)}{\partial x},\;\frac{\partial(yz)}{\partial x}-\frac{\partial(x^2)}{\partial y}\right)$$
$$= (0 - y,\;\; 0 - z,\;\; 0 - 0) = \boxed{(-y,\;-z,\;0)}$$

\textbf{Vérification :} $\text{rot}\,\mathbf{F}\neq\mathbf{0}$ donc $\mathbf{F}$ n'est pas un champ gradient.
\end{reponsebox}

% ═══════════════════════════════════════════════════════════════════════════════
\section{Partie 7 - Applications géométriques}
% ═══════════════════════════════════════════════════════════════════════════════

%────── Q26 ────────────────────────────────────────────────────────────────────
\qtitre{QCM}{C - $\nabla F(M)$}
\qenonce{Le vecteur normal à la surface $F(x,y,z)=0$ en un point $M$ est donné par :
\quad A)~$\text{div}\,F(M)$
\quad B)~$\Delta F(M)$
\quad C)~$\nabla F(M)$
\quad D)~$\text{rot}\,\nabla F(M)$}

\begin{methodebox}
\textbf{Vecteur normal à une surface implicite.} La surface $\mathcal{S} : F(x,y,z)=0$ est une surface de niveau de $F$. Le gradient $\nabla F$ est perpendiculaire aux surfaces de niveau, donc $\nabla F(M)$ est le \textbf{vecteur normal} à $\mathcal{S}$ en $M$.

La proposition D est nulle (rot grad = 0 par Schwartz).
\end{methodebox}

%────── Q27 ────────────────────────────────────────────────────────────────────
\qtitre{QCM}{B - $\frac{\partial F}{\partial x}(M_0)(x-x_0)+\frac{\partial F}{\partial y}(M_0)(y-y_0)+\frac{\partial F}{\partial z}(M_0)(z-z_0)=0$}
\qenonce{L'équation du plan tangent à $F(x,y,z)=0$ au point $M_0=(x_0,y_0,z_0)$ est :
\quad A)~$F(x,y,z)=F(x_0,y_0,z_0)$
\quad B)~$\frac{\partial F}{\partial x}(M_0)(x-x_0)+\frac{\partial F}{\partial y}(M_0)(y-y_0)+\frac{\partial F}{\partial z}(M_0)(z-z_0)=0$
\quad C)~$\nabla F(M_0)\cdot(x,y,z)=0$
\quad D)~$F(x-x_0,y-y_0,z-z_0)=0$}

\begin{methodebox}
\textbf{Plan tangent à une surface implicite.} Le plan tangent en $M_0$ a pour vecteur normal $\nabla F(M_0)$ et passe par $M_0$ :
$$\nabla F(M_0)\cdot\overrightarrow{M_0M} = 0 \quad\Longleftrightarrow\quad \frac{\partial F}{\partial x}(x-x_0)+\frac{\partial F}{\partial y}(y-y_0)+\frac{\partial F}{\partial z}(z-z_0) = 0$$
La proposition A est vraie uniquement si $F(M_0)=0$ et ne définit pas un plan. La proposition C oublie que le plan passe par $M_0$ (manque la translation).
\end{methodebox}

%────── Q28 (Calcul) ──────────────────────────────────────────────────────────
\qtitre{Calcul}{$\nabla F(1,2,2) = (2,4,4)$ \quad et \quad plan tangent : $x+2y+2z=9$}
\qenonce{Soit la sphère $F(x,y,z)=x^2+y^2+z^2-9$. Calculer $\nabla F$ en $M_0=(1,2,2)$ et donner l'équation du plan tangent.}

\begin{reponsebox}
\textbf{Calcul du gradient :}
$$\nabla F = \left(\frac{\partial F}{\partial x},\;\frac{\partial F}{\partial y},\;\frac{\partial F}{\partial z}\right) = (2x,\;2y,\;2z)$$
$$\boxed{\nabla F(1,2,2) = (2,\;4,\;4)}$$

\textbf{Vérification que $M_0$ est sur la surface :} $1^2+2^2+2^2 = 1+4+4 = 9$ \checkmark

\textbf{Équation du plan tangent :}
$$2(x-1) + 4(y-2) + 4(z-2) = 0$$
$$2x - 2 + 4y - 8 + 4z - 8 = 0$$
$$\boxed{x + 2y + 2z = 9}$$

\textbf{Interprétation géométrique :} le plan tangent à la sphère de rayon 3 en $(1,2,2)$ est perpendiculaire au rayon $\overrightarrow{OM_0}=(1,2,2)$, vecteur directeur normal du plan.
\end{reponsebox}

% ═══════════════════════════════════════════════════════════════════════════════
\newpage
\section*{\color{bleutexte}Synthèse des méthodes à retenir}
% ═══════════════════════════════════════════════════════════════════════════════

\begin{methodebox}
\textbf{Calcul différentiel - Procédure}
\begin{enumerate}[itemsep=2pt, topsep=2pt]
  \item Différentielle : $df = \dfrac{\partial f}{\partial x}dx + \dfrac{\partial f}{\partial y}dy$ (généralise la dérivée à plusieurs variables)
  \item Plan tangent à $z=f(x,y)$ en $(a,b)$ : $z = f(a,b) + \dfrac{\partial f}{\partial x}(a,b)(x-a) + \dfrac{\partial f}{\partial y}(a,b)(y-b)$
  \item Règle de la chaîne : $\dfrac{\partial(g\circ h)}{\partial x} = g'(h)\cdot\dfrac{\partial h}{\partial x}$
  \item Schwartz ($f\in\mathcal{C}^2$) : $\dfrac{\partial^2 f}{\partial x\partial y} = \dfrac{\partial^2 f}{\partial y\partial x}$
\end{enumerate}
\end{methodebox}

\vspace{6pt}

\begin{methodebox}
\textbf{Extrema - Classification des points critiques}

En un point critique (où $\nabla f = \mathbf{0}$), calculer $H = f_{xx}\,f_{yy} - f_{xy}^2$ :
\begin{center}
\begin{tabular}{@{}lll@{}}
\toprule
Condition & Nature & \\
\midrule
$H > 0$ et $f_{xx} > 0$ & Minimum local    & $f$ croît dans toutes les directions \\
$H > 0$ et $f_{xx} < 0$ & Maximum local    & $f$ décroît dans toutes les directions \\
$H < 0$                 & Point selle (col) & croissance dans certaines directions \\
$H = 0$                 & Cas douteux      & critère insuffisant \\
\bottomrule
\end{tabular}
\end{center}
\end{methodebox}

\vspace{6pt}

\begin{methodebox}
\textbf{Opérateurs vectoriels - Récapitulatif}
\begin{center}
\begin{tabular}{@{}llll@{}}
\toprule
Opérateur & Notation & Formule & Nature du résultat \\
\midrule
Gradient   & $\nabla f$           & $\bigl(\partial_x f,\,\partial_y f,\,\partial_z f\bigr)$                                         & Vecteur \\[4pt]
Divergence & $\text{div}\,\mathbf{F}$ & $\partial_x P + \partial_y Q + \partial_z R$                                                  & Scalaire \\[4pt]
Rotationnel & $\text{rot}\,\mathbf{F}$ & $\bigl(\partial_y R - \partial_z Q,\;\partial_z P - \partial_x R,\;\partial_x Q - \partial_y P\bigr)$ & Vecteur \\[4pt]
Laplacien  & $\Delta f$           & $\partial^2_{xx} f + \partial^2_{yy} f + \partial^2_{zz} f$                                    & Scalaire \\
\bottomrule
\end{tabular}
\end{center}

\textbf{Identités fondamentales :} $\quad\text{rot}(\nabla f) = \mathbf{0}$ \quad et \quad $\text{div}(\text{rot}\,\mathbf{F}) = 0$
\end{methodebox}

\vspace{6pt}

\begin{methodebox}
\textbf{Plan tangent à une surface implicite $F(x,y,z)=0$}

\textbf{Étape 1 :} Calculer $\nabla F = \bigl(\dfrac{\partial F}{\partial x},\,\dfrac{\partial F}{\partial y},\,\dfrac{\partial F}{\partial z}\bigr)$.

\textbf{Étape 2 :} Évaluer en $M_0$ pour obtenir le vecteur normal $\mathbf{n} = \nabla F(M_0)$.

\textbf{Étape 3 :} Équation du plan tangent : $\mathbf{n}\cdot\overrightarrow{M_0M} = 0$, soit :
$$\frac{\partial F}{\partial x}(M_0)(x-x_0) + \frac{\partial F}{\partial y}(M_0)(y-y_0) + \frac{\partial F}{\partial z}(M_0)(z-z_0) = 0$$

\textbf{Vérification :} s'assurer que $M_0$ vérifie bien $F(M_0) = 0$ avant de commencer.
\end{methodebox}

\end{document}
